
\setcounter{subsection}{4}
\subsection{}
\label{XIII.4.5}
Retomemos las notaciones y las hipótesis de~\Ref{XIII.4.3}. Si
$\bar{s}$ es un punto geométrico de~$S$ y $a=g(\bar{s})$,
la sección~$g$ permite definir un morfismo
$$
w\colon \pi_1(S,a)\to \pi_1'(X,a),
$$
de modo que $\pi_1'(X,a)$ se identifica con el producto semidirecto de
$\pi_1(S,a)$ por
\ifthenelse{\boolean{orig}}
{$\pi_1(X_{\bar{s}},a)$}
{$\pi_1^\LL(X_{\bar{s}},a)$}.
El grupo profinito $\pi_1(S,a)$ actúa, por tanto, sobre
\ifthenelse{\boolean{orig}}
{$\pi_1(X_{\bar{s}},a)$}
{$\pi_1^\LL(X_{\bar{s}},a)$}.
Como $\pi_1^\LL(X_{\bar{s}},a)$ es límite proyectivo estricto de grupos
invariantes bajo la acción de~$\pi_1(S,a)$, dar
$\pi_1^\LL(X_{\bar{s}},a)$ provisto de esta acción equivale
a dar un sistema proyectivo estricto de esquemas
en grupos finitos y étales sobre~$S$, que se denota por
$$
\pi_1^\LL(X/S,g,\bar{s})\quad\text{o
simplemente}\quad\pi_1^\LL(X/S,g).
$$
\label{indnot:mk}\oldindexnot{$\pi_1^\LL(X/S,g,\bar{s})$ o $\pi_1^\LL(X/S,g)$|hyperpage}%

Se tienen entonces las propiedades siguientes:
\setcounter{subsubsection}{0}

\subsubsection{}
\label{XIII.4.5.1}
Para todo esquema en grupos finito y étale~$G$ sobre~$S$, cuyas
fibras son~$\LL$-grupos, el conjunto~$E$ de las clases de torsores
$P$ bajo la imagen inversa~$G_X$ de~$G$ sobre~$X$, provistos de un
isomorfismo $g^*P\isomto G$, es canónicamente isomorfo al conjunto
$$
\Hom_S(\pi_1^\LL(X/S,g,\bar{s}),G) \quad\text{módulo automorfismos
interiores de~$G$.}
$$

\subsubsection{}
\label{XIII.4.5.2}
Para todo esquema en grupos finito y étale~$G$ sobre~$S$ cuyas
fibras son~$\LL$-grupos, el haz $\R^1f_*G_X$ es
canónicamente isomorfo al haz asociado al prehaz
$$
S'\mto \Hom_{S'}(\pi_1^\LL(X/S,g,\bar{s}),G)\
\text{módulo automorfismos interiores de~$G$.}
$$
($S'$ designa un esquema étale sobre~$S$).

\subsubsection{}
\label{XIII.4.5.3}
Sean
\marginpar{421}
$S'$ un $S$-esquema conexo, $\bar{s}$ un punto geométrico de~$S'$, y~$X'$,
$g'$ las imágenes inversas respectivas de~$X,g$ sobre~$S'$. Entonces,
$\pi_1^\LL(X'/S',g',\bar{s})$ es canónicamente isomorfo a
la imagen inversa de~$\pi_1^\LL(X/S,g,\bar{s})$ sobre~$S'$. Para todo
punto geométrico~$\xi$ de~$S$, la fibra
$\pi_1^\LL(X/S,g,\bar{s})_\xi$ es isomorfa a~$\pi_1^\LL(X_\xi)$.

Dar~$G$ equivale, en efecto, a dar un~$\LL$-grupo abstracto~$\mathbf{G}$
sobre el cual actúa~$\pi_1(S,a)$, de donde resulta una acción de
$\pi_1'(X,a)$ sobre~$\mathbf{G}$. El isomorfismo definido
en~\Ref{XIII.4.5.1} se obtiene entonces, por restricción al subconjunto
$E$, a partir del morfismo canónico
\[
\H^1(\pi_1'(X,a),\mathbf{G})\to
\H^1(\pi_1^\LL(X_{\bar{s}},a),\mathbf{G}) =
\Hom(\pi_1^\LL(X_{\bar{s}},a),\mathbf{G})/
\ifthenelse{\boolean{orig}}
{\text{aut. int. }G}
{\text{aut. int. }\mathbf{G}},
\]
y este morfismo envía biyectivamente~$E$ sobre el subconjunto de los
morfismos de~$\pi_1^\LL(X_{\bar{s}},a)$ a~$\mathbf{G}$ compatibles con
la acción de~$\pi_1(S,a)$. La afirmación~\Ref{XIII.4.5.3} resulta
entonces de la definición de~$\pi_1^\LL(X/S,g,\bar{s})$, habida cuenta
de la sucesión exacta de homotopía~\eqref{eq:XIII.4.3.1}, y
\Ref{XIII.4.5.2} se deduce de~\Ref{XIII.4.5.1} y~\Ref{XIII.4.5.3}.
